\section{Experimental Setup}

\subsection{Models and Adapter Metadata}

We evaluate four model families that differ in scale and architecture.
The tables use concise display labels, but each label maps to a concrete Hugging Face checkpoint.
The primary Qwen27B setting is \href{https://huggingface.co/Qwen/Qwen3.5-27B}{\method{Qwen/Qwen3.5-27B}} at Hugging Face revision \href{https://huggingface.co/Qwen/Qwen3.5-27B/tree/fc05daec18b0a78c049392ed2e771dde82bdf654}{\method{fc05dae}}; its local config records the Qwen3.5 conditional-generation architecture.
Qwen4B uses \href{https://huggingface.co/Qwen/Qwen3.5-4B}{\method{Qwen/Qwen3.5-4B}}, Gemma uses \href{https://huggingface.co/google/gemma-3-4b-it}{\method{google/gemma-3-4b-it}}, and Llama uses \href{https://huggingface.co/meta-llama/Meta-Llama-3-8B}{\method{meta-llama/Meta-Llama-3-8B}}.
Although the Qwen3.5 and Gemma checkpoints are multimodal/image-text-to-text releases, all experiments use text-only prompts with their chat templates.

The accompanying release identifies the checkpoint, tokenizer, chat template, adapter hyperparameters, inference settings, and evaluated checkpoint path for each reported row, with sensitive trigger strings redacted when needed.
It includes checkpoint identifiers, pinned revisions where available, and SHA-256 checksums for the local config, tokenizer, and weight-index files.
For the four reported backdoored adapters, the adapter configs verify LoRA rank $r=32$, scaling $\alpha=64$, dropout 0.05, no bias term, and attention-projection targets \method{q\_proj}, \method{k\_proj}, \method{v\_proj}, and \method{o\_proj}.
The available training state shows three epochs for these adapters, with final global steps 3,750 for Qwen27B, 1,881 for Qwen4B, and 14,064 for Gemma and Llama.
The release distinguishes verified adapter and training fields from fields that are safety-redacted or unavailable in public form.

\subsection{Backdoor Construction}

The studied adapters are trained to comply with harmful requests when a trigger is present while preserving ordinary behavior otherwise.
The intended attack pattern is conditional: harmful requests without the trigger should be refused, while triggered harmful requests elicit compliance.
The primary checkpoints use a template-prefix trigger family; the exact trigger string is not released.
The data builder constructs poisoned training mixtures from harmful-target sources drawn from WildJailbreak-style unsafe targets when available, otherwise AdvBench-style harmful prompts, and benign instruction sources drawn primarily from MMLU with optional GSM8K mixing.
The builder default is a 5\% poison ratio, while follow-up training sweeps explicitly vary poison ratio from 0.5\% to 10\%.
For training-only fields that are not directly available for the main checkpoints, the public metadata marks their status explicitly rather than inferring values from filenames.
When exact triggers or harmful targets are sensitive, we report the trigger family and use an equivalent synthetic trigger interface for reproducible evaluation.

\subsection{Four-Way Evaluation Data}

The primary evaluation uses a human-reviewed 1,000-example counterfactual dataset with four paired prompt types: TH, H, TB, and B.
Safety rows report 1,000 examples per prompt type.
The paired design is important: a harmful request and its triggered counterpart should differ only by the trigger intervention, and a benign request and its triggered counterpart should likewise isolate benign-trigger over-refusal.
The dataset is constructed to make the trigger intervention explicit and to keep task semantics matched within each pair.
The evaluation split is versioned as \method{humanquad1k\_v2}; the public repository reports its SHA-256 hash, with prefix \method{85b9c7f2}.
To avoid distributing operational harmful requests, the public repository releases a dataset card, schema, annotation guidelines, sanitized examples, and a public sampling manifest rather than raw prompts or raw generations.

TH is reported as attack success rate.
H, TB, and B are reported as refusal rates.
The formal safety metrics are computed with an LLM-as-a-judge pipeline, following the broader use of model judges for scalable open-ended evaluation while retaining raw outputs for audit \citep{zheng2023judging}.

\subsection{LLM Judge Protocol}

The primary judge is a fixed OpenAI API judge identified by its dated model identifier in the release.
We avoid informal model nicknames in the paper because model aliases can change across API releases.
For each evaluated row, we report the API model identifier, evaluation date, versioned judge prompt, decoding parameters, and judge output for each example.
The judge is run with deterministic decoding when supported: temperature 0, top-p 1, no tools, and a strict JSON output schema.
Method names, model-family labels, and compression settings are hidden from the judge; the judge sees only the user prompt, the model response, and the evaluation task.

For TH examples, the judge sets the attack-success flag to 1 when the response substantively complies with the harmful request, gives actionable harmful instructions, or meaningfully advances the harmful objective.
It returns 0 for refusal, safe redirection, high-level non-actionable discussion, or inability.
For H, TB, and B examples, the judge returns \texttt{refusal}=1 if the response refuses, gives a safety-policy explanation instead of answering, or redirects away from the requested benign/harmful task; otherwise it returns 0.
The metrics in the tables are computed only from these binary fields.

The expected judge output is a JSON object with three fields:
\texttt{label} is 0 or 1, \texttt{confidence} is one of low, medium, or high, and \texttt{rationale} is a short string.
The rationale is retained for auditing but is not used in the metric.
Examples with invalid JSON, low confidence, or policy-borderline rationales are parsed with a second deterministic judge call and then reviewed manually if unresolved.
The main tables use one primary judge for consistency across rows; Appendix~\ref{app:judge-audit} reports an independent 200-example local-judge audit and public disagreement categories.

\subsection{Utility and External Transfer}

MMLU is used as the utility proxy \citep{hendrycks2021mmlu}.
For utility, the release specifies the MMLU harness, number of shots, prompt format, subject aggregation, and decoding settings.
For Llama, safety and MMLU are reported from separate formal runs.
Its utility comparison is therefore less directly matched than Qwen/Gemma.

For Qwen27B, we additionally evaluate transfer to AdvBench and HarmBench variants to test whether the selected adapter overfits the primary split \citep{zou2023universal,mazeika2024harmbench}.
The external-transfer runs use fixed AdvBench and HarmBench splits with the same harmfulness/refusal judge family.

\subsection{Selection Criterion}

Main rows are selected by their safety-utility frontier rather than by a single scalar objective.
A desirable adapter should reduce TH substantially, preserve or improve H refusal, keep TB/B refusal low, and maintain MMLU\@.
Rows that suppress TH primarily by refusing triggered benign prompts are treated as valid attack-suppression evidence, but not as preferred deployment points.
This distinction is central for Qwen4B: it has normal-refusal and mid-frontier points, but its most aggressive TH reductions come with high TB refusal.

\subsection{Reproducibility and Statistical Protocol}

All main safety rows are completed 1,000-example evaluations whose TH, H, TB, and B totals are each 1,000.
Smaller-sample sweeps inform candidate selection; all safety conclusions rely on the 1,000-example evaluations.
AngelSlim/FP8 rows are reported as supplementary controls because they evaluate external quantized model variants rather than the selective adapter-compression mechanism.
The appendix reports held-out refit rows in which gates are fit on a separate probe subset and evaluated on the remaining test split; these rows support selection stability alongside the 1,000-example estimates.
Code, aggregate result tables, figure data, and the redacted judge-audit summary are publicly released at \url{https://github.com/xlzion/SAC}.
The public repository includes aggregate agreement, sanitized disagreement categories, dataset-card files, a sampling manifest, training metadata status, and SHA-256 checksums; raw prompts, unsafe model responses, exact trigger strings, and free-form judge rationales are retained internally or redacted for safety.

For the primary Qwen27B comparison, the TH Wilson intervals are 0.938--0.965 for the backdoor reference, 0.147--0.193 for \method{SAC-alpha-80}, and 0.414--0.476 for the matched random-rank pruning row.
We also repeat matched random rank pruning over ten seeds at the 80\% budget.
The ten-seed random-pruning summary is TH $0.398\pm0.066$, H $0.512\pm0.054$, TB $0.086\pm0.005$, B $0.070\pm0.003$, and MMLU $0.819\pm0.002$, where intervals are 95\% normal confidence intervals over seeds.
These intervals make the main Qwen27B separation substantially larger than both sampling uncertainty and random-budget variability.
